\textbf{Figure 4. Single-cell programs in held-out livers and three results inside hepatocytes.}

Program memberships and weights remain fixed. Single-cell donor analyses use deduplicated unique physical cells (650,083 unique cells across 145 captures and 104 source donors; 226,457 unsorted unique cells across 63 donors), with source donor identifiers as the inferential unit; bulk analyses use selected genetic-identity clusters. All associations are cross-sectional.

\textbf{(A)} Seven datasets contribute 650,083 unique cells across 145 physical captures and 104 source donors (historical run-level discovery inputs held 1,232,318 cells across 275 libraries). The 117 programs were discovered in five lineage subsets and remain frozen. Source corrections exclude GSE202379 P98 from disease inference because its control diagnosis conflicts with SAF S1A1F1, retain P30 as a qualified nonsteatotic reference (S0A1F0), and place GSE244832 D11/D14 in MASL/MASH per deposited RNA metadata. RNA diagnoses do not verify RNA--ATAC pair identities. Histology, diagnosis and etiology remain separate.

\textbf{(B)} The unchanged uniform 50,000-cell UMAP sample (seed 42), colored by cell type. Cells are descriptive observations, not independent replicates. Disease-density overlays are omitted because the diagnosis groups contain different dataset mixtures. Colors shared with E denote discovery lineage.

\textbf{(C)} Discovery versus content-associated lineage for all 117 programs on deduplicated unsorted donor pseudobulk across six active lineages (excluding sorted GSE136103 cells): 78 match (66.7\%, 95\% donor-bootstrap CI 65.0--67.5\%) and 39 differ. Gray cells mark the diagonal. Content association does not establish transcript origin.

\textbf{(D)} Each point is one of 112 testable frozen Hotspot programs in SteatoSITE bulk RNA (493 staged donors; F0--F1/F2/F3/F4 = 255/74/75/89). x, unadjusted program score change per stage step (\texttt{score\ \textasciitilde{}\ stage\ +\ batch\ +\ inferred\_\allowbreak{}sex}, HC3; score SD); y, the same coefficient adjusted for seven isometric log-ratio coordinates from 22-autosome leave-one-chromosome-out (LOCO) InstaPrism deconvolution separating hepatocytes and cholangiocytes (vertical bars span the 22-autosome LOCO range around the modal-chromosome estimate). Dashed line, y = x. Colour, unadjusted BH q \textless{} 0.05 over 112 programs: 63 up, 35 down, 14 not significant; 4 programs retain LOCO BH q \textless{} 0.05 across all 22 chromosomes. Labeled: Stromal ECM survives 22/22 LOCO chromosomes (modal chr3 $\beta$ = 0.0881, q = 8.93 $\times$ 10$^{-7}$), whereas Ductular injury attenuates under modal-chr2 LOCO ($\beta$ = 0.0297, q = 0.402; 0/22 chromosomes q \textless{} 0.05). Also labeled: Complement, Nucleoside catabolism and NRF2 antioxidant. Cross-sectional stage-associated remodeling; unit, donor.

\textbf{(E)} GNMT, MAT1A and CYP2C19 in hepatocyte pseudobulk, steatohepatitis or worse minus control (log2 CPM, 95\% CI; HC3-studentized Freedman--Lane permutations). Labels give control/disease donors. Circles, cohorts where the decrease was first seen (GSE298719, Flex probe panel without CYP2C19; GSE256398); triangles, held-out cohorts. Model in every cohort: \texttt{logCPM\ \textasciitilde{}\ disease\ +\ sex\ +\ log10(hepatocyte\ nuclei)}. GSE289173 (23 control, 17 steatohepatitis donors after excluding GSM8785704; controls collected as a separate series): CYP2C19 -1.69 (-2.56 to -0.81, Holm \emph{p} = 0.0034), GNMT -1.17 (-1.92 to -0.41, Holm \emph{p} = 0.0045), MAT1A -0.92 (-1.44 to -0.41, Holm \emph{p} = 0.0034). GSE202379 and GSE244832 combined (8 healthy, 44 disease deduplicated donors): uninformative (GNMT Holm \emph{p} = 0.750, MAT1A \emph{p} = 0.691, CYP2C19 \emph{p} = 0.691). Unit, donor. The same model for the eight recurring example genes is in Figure S4AG.

\textbf{(F)} Fixed 64-gene NF-$\kappa$B set (defined in GSE298719 and GSE256398; 62 measured) in CCL20+ hepatocytes minus depth- and injury-matched CCL20$-$ hepatocytes of the same deduplicated donor, minus the mean of 20 random anchor genes (log2). One point per donor with at least 10 CCL20+ hepatocytes; bar, median. GSE244832, 8 deduplicated donors: mean +0.85, 8 of 8 above zero, exact sign-flip \emph{p} = 0.0078, Holm \emph{p} = 0.0156 over two cohorts. GSE202379, 5 donors: mean +1.01, 5 of 5 above zero, \emph{p} = 0.0625, the smallest value attainable with 5 donors. Exploratory; the set was fixed before these cohorts were read.

\textbf{(G)} Share of each gene's UMIs from cells labeled hepatocytes, per deduplicated MASLD donor (bar, median over donors with at least one UMI) in GSE298719, GSE244832 and GSE202379. Gray, reference genes ALB and PTPRC. Medians for CCL20: 0.75 (GSE298719), 0.88 (GSE244832) and 0.50 (GSE202379) raw (0.90 in GSE256398), and 0.34, 0.61 and 0.33 (0.61 in GSE256398) after strict PTPRC and ductular-marker removal. Donors with zero UMIs are not plotted. In the nuclei cohorts PTPRC sets a non-hepatocyte floor from ambient RNA, so raw shares there are upper bounds. Descriptive; no test.

\textbf{(H)} Change in a five-gene bulk inflamed-hepatocyte score (FCAMR, TRIM31, CXCL10, IL32, CCL20; within-cohort z of log2 TMM CPM) per PNPLA3 rs738409 148M allele called from RNA reads (\texttt{score\ \textasciitilde{}\ allele\ +\ cohort\ +\ sex\ +\ 4\ ancestry\ PCs}, HC3, 95\% CI), five bulk cohorts. All + 4 PCs: +0.153 SD (0.070 to 0.236; \emph{p} = 3.26 $\times$ 10$^{-4}$; n = 757); EUR + 4 PCs: +0.147 SD (0.059 to 0.236; n = 641); + Clean histology + 4 PCs (excluding unstated GSE213621 Control fibrosis codes): +0.029 SD (-0.045 to 0.102; \emph{p} = 0.441; n = 633); EUR + clean histology + 4 PCs: +0.026 SD (n = 547; \emph{p} = 0.510); MASLD only + 4 PCs: +0.056 SD (n = 626; \emph{p} = 0.184). Note: GSE135251 rs738409 HWE \emph{p} = 0.018 (EUR 0.032). Exploratory.

See also Figure S4 and Supplementary Notes 8 and 11.
